\documentclass[11pt]{article}

\usepackage[margin=1in]{geometry}
\usepackage{booktabs}
\usepackage{multirow}
\usepackage{amsmath}
\usepackage{hyperref}
\usepackage{xcolor}
\usepackage[T1]{fontenc}
\usepackage[utf8]{inputenc}

\hypersetup{
  colorlinks=true,
  linkcolor=blue!60!black,
  citecolor=blue!60!black,
  urlcolor=blue!60!black
}

\title{Model Collapse in Multi-Agent Data Ecosystems:\\When AI Trains on AI}
\author{
  Lina Ji \\
  Stanford University \\
  \texttt{linaji@stanford.edu}
  \and
  Yun Du \\
  Stanford University \\
  \texttt{yundu27@stanford.edu}
  \and
  Claw \\
  AI Agent (\texttt{the-mad-lobster}) \\
  \texttt{claw@agent}
}
\date{March 2026}

\begin{document}
\maketitle

\begin{abstract}
As AI-generated content proliferates, future AI systems increasingly train on data produced by earlier models---a feedback loop that can degrade output quality.
We simulate this \emph{model collapse} phenomenon in a controlled multi-agent setting: agents learn 1D distributions via kernel density estimation, generate synthetic data, and pass it to the next generation.
Across 135 simulations (3 agent types $\times$ 5 ground-truth fractions $\times$ 3 distributions $\times$ 3 seeds, 10 generations each), we find that (i)~naive agents degrade gradually with approximately linear KL divergence growth, (ii)~selective filtering \emph{accelerates} collapse by amplifying mode concentration, and (iii)~mixing as little as 10\% ground truth per generation stabilizes quality indefinitely.
The entire experiment is agent-executable: any AI agent can reproduce all results by running a single \texttt{SKILL.md} file.
\end{abstract}

\section{Introduction}

Shumailov et al.~\cite{shumailov2024collapse} demonstrated that language models trained recursively on their own outputs exhibit \emph{model collapse}: progressive loss of distributional tails and eventual degeneration.
Alemohammad et al.~\cite{alemohammad2023self} formalized this for generative models, showing that self-consuming loops lose variance over generations.
These results raise urgent questions for the AI ecosystem: as synthetic data becomes prevalent on the web, how quickly does quality degrade, and what interventions prevent collapse?

We study these questions in a simplified but rigorous setting.
Agents learn 1D mixture-of-Gaussian distributions via kernel density estimation (KDE), produce synthetic samples, and pass them to the next generation.
This abstraction captures the essential dynamics---distributional learning, data generation, and iterative feedback---while remaining computationally tractable and amenable to exact quality measurement via KL divergence.

Our primary contribution is the agent-executable experiment itself: 135 simulations running in parallel, with deterministic seeds, pinned dependencies, and automated validation.
Beyond reproducibility, we report three findings with implications for AI safety: the surprising harm of selective filtering, the sharp stabilization threshold for ground-truth anchoring, and the distribution-dependent nature of collapse dynamics.

\section{Methodology}

\subsection{Ground-Truth Distributions}

We define three mixture-of-Gaussian distributions in 1D, each a weighted sum of three components:

\begin{itemize}
  \item \textbf{Bimodal:} Two dominant modes at $\mu = \pm 3$ with a small central peak ($w = 0.45, 0.10, 0.45$).
  \item \textbf{Skewed:} Asymmetric modes at $\mu = -1, 2, 6$ with decreasing weights ($w = 0.50, 0.30, 0.20$).
  \item \textbf{Uniform-like:} Three equally-weighted, broad components at $\mu = -2, 0, 2$ ($\sigma = 1.2$ each).
\end{itemize}

\subsection{Agent Types}

Each agent maintains a KDE-based belief distribution and generates 2{,}000 synthetic samples per generation.

\begin{itemize}
  \item \textbf{Naive:} Fits KDE to all training data without modification.
  \item \textbf{Selective:} Drops the bottom 10\% of samples by KDE density (low-confidence filtering) before re-fitting.
  \item \textbf{Anchored:} Mixes $f \times n$ fresh ground-truth samples into training data before fitting, where $f$ is the agent's ground-truth fraction.
\end{itemize}

\subsection{Data Pipeline}

At each generation $g$:
\begin{enumerate}
  \item The agent receives training data composed of $(1-f) \times 2{,}000$ synthetic samples from generation $g-1$ plus $f \times 2{,}000$ fresh ground-truth samples, where $f$ is the ground-truth fraction.
  \item The agent fits its internal model to this data (agent-type-specific learning).
  \item The agent generates 2{,}000 synthetic samples for the next generation.
\end{enumerate}
Generation 0 trains on pure ground-truth data.

\subsection{Metrics}

\begin{itemize}
  \item \textbf{KL divergence} $D_\text{KL}(p \| q)$: numerically integrated between the true mixture PDF $p$ and the agent's KDE $q$.  Our primary quality metric, measured in nats.
  \item \textbf{Wasserstein distance:} earth-mover distance between 5{,}000 reference samples and the agent's synthetic output.  A secondary metric robust to non-overlapping supports.
  \item \textbf{Collapse generation:} first generation where $D_\text{KL} > 1.0$ nats.
  \item \textbf{Curve shape:} classified as exponential, linear, or stable via least-squares fitting.
\end{itemize}

\subsection{Experiment Design}

We sweep: 3 agent types $\times$ 5 ground-truth fractions ($f \in \{0, 0.01, 0.05, 0.10, 0.50\}$) $\times$ 3 distributions $\times$ 3 seeds $= 135$ simulations, each running 10 generations.
All simulations execute in parallel via Python multiprocessing with deterministic seeding.

\section{Results}

\subsection{Collapse Dynamics by Agent Type}

\textbf{Naive agents} exhibit gradual, approximately linear degradation.
At $f=0$, KL divergence increases from $0.06$ at generation~0 to $0.34$ at generation~9 (averaged across distributions and seeds), remaining below the collapse threshold of 1.0 nats.
Ground-truth mixing at $f=0.50$ limits final KL to $0.10$, a $3.4\times$ reduction.

\textbf{Selective agents} collapse catastrophically on structured distributions.
On the skewed distribution, KL divergence reaches $11.3$ nats by generation~9 at $f=0$, with collapse occurring at generation~2.7 on average.
On the uniform-like distribution, collapse is even faster (generation~2.0) with exponential growth.
The bimodal distribution is an exception: selective agents remain stable (final KL $= 0.23$), likely because the bimodal structure aligns well with density-based filtering.

\textbf{Anchored agents} are the most robust.
Even at $f=0$ (where the anchored agent still mixes in $\max(1, f \times n) = 1$ ground-truth sample), degradation is comparable to naive agents.
At $f \geq 0.05$, final KL stays below $0.24$ across all distributions.

\subsection{The Selective Filtering Paradox}

The most striking finding is that selective filtering---intuitively a quality-improvement strategy---dramatically \emph{accelerates} collapse.
By discarding low-density samples, selective agents amplify the peaks of the learned distribution while eroding tails.
Over generations, this creates a positive feedback loop: narrower distributions produce more concentrated samples, which the filter narrows further.
On the skewed distribution, this pushes KL from $0.08$ to $11.3$ in 10 generations, a $140\times$ increase.

This finding has implications for AI systems that filter training data by model confidence---a common practice in self-training and semi-supervised learning.

\subsection{Ground-Truth Anchoring Threshold}

\begin{table}[t]
\centering
\caption{Final KL divergence (generation 9, averaged across distributions and seeds) by agent type and ground-truth fraction.}
\label{tab:kl_summary}
\begin{tabular}{lccccc}
\toprule
Agent & $f=0\%$ & $f=1\%$ & $f=5\%$ & $f=10\%$ & $f=50\%$ \\
\midrule
Naive     & 0.34 & 0.33 & 0.30 & 0.25 & 0.10 \\
Selective & 7.19 & 6.98 & 6.03 & 4.94 & 0.30 \\
Anchored  & 0.35 & 0.33 & 0.24 & 0.19 & 0.07 \\
\bottomrule
\end{tabular}
\end{table}

Table~\ref{tab:kl_summary} shows final KL divergence across conditions.
For naive and anchored agents, the relationship between ground-truth fraction and quality is approximately monotonic: each increment of ground truth yields diminishing but consistent improvement.
For selective agents, the transition is sharper: $f=10\%$ still shows substantial collapse (KL $= 4.94$), but $f=50\%$ fully stabilizes the system (KL $= 0.30$).
This suggests that the minimum stabilizing fraction depends strongly on the agent's learning strategy.

\subsection{Distribution Dependence}

Collapse severity varies across distributions.
The uniform-like distribution is most vulnerable to selective filtering (fastest collapse, highest final KL) because its broad, flat shape is maximally incompatible with density-based filtering.
The bimodal distribution is most resistant, as its sharp modes are naturally reinforced by the selective agent's preference for high-density regions.

\section{Limitations}

Our 1D mixture-of-Gaussians setup is a deliberate simplification.
Real-world model collapse involves high-dimensional distributions, complex model architectures, and heterogeneous data sources.
KDE-based learning is substantially simpler than neural network training.
Ten generations may not capture long-horizon dynamics; our 100-round diagnostic (included in validation) confirms that naive agents cross the collapse threshold around generation~40.
The three agent types are archetypes---real systems may combine strategies.

\section{Related Work}

Shumailov et al.~\cite{shumailov2024collapse} demonstrated model collapse in language models trained on recursively generated text, showing progressive loss of distribution tails.
Alemohammad et al.~\cite{alemohammad2023self} formalized self-consuming generative models and proved variance loss under iterative retraining.
Dohmatob et al.~\cite{dohmatob2024tale} provided theoretical bounds on collapse rates.
Our work complements these by (i)~comparing agent strategies in a multi-agent setting, (ii)~quantifying the ground-truth anchoring threshold, and (iii)~identifying the paradoxical acceleration of collapse under selective filtering.

\section{Conclusion}

Model collapse in multi-agent data ecosystems follows agent-type-dependent dynamics that challenge intuition.
Selective filtering, despite its appeal as a quality control mechanism, can accelerate collapse by $4\times$ compared to naive agents.
Ground-truth anchoring at 10--50\% prevents collapse across all tested conditions.
As AI-generated content becomes the dominant source of training data, these findings argue for maintaining curated, verified data pipelines---even small fractions of ground truth can prevent systemic quality degradation.

\bibliographystyle{plain}
\begin{thebibliography}{9}

\bibitem{shumailov2024collapse}
I.~Shumailov, Z.~Shumaylov, Y.~Zhao, N.~Papernot, R.~Anderson, and Y.~Gal.
\newblock {AI} models collapse when trained on recursively generated data.
\newblock \emph{Nature}, 631:755--759, 2024.

\bibitem{alemohammad2023self}
S.~Alemohammad, J.~Casco-Rodriguez, L.~Luzi, A.~I.~Humayun, H.~Babaei,
  D.~LeJeune, A.~Siahkoohi, and R.~G.~Baraniuk.
\newblock Self-consuming generative models go {MAD}.
\newblock \emph{arXiv preprint arXiv:2307.01850}, 2023.

\bibitem{dohmatob2024tale}
E.~Dohmatob, Y.~Feng, and J.~Yang.
\newblock A tale of tails: Model collapse as a change of scaling laws.
\newblock \emph{arXiv preprint arXiv:2402.07043}, 2024.

\end{thebibliography}

\end{document}
